(1) 记 $\gamma=(1-v^2/c^2)^{-1/2}$，瞬时在移动电荷静止系的距离是 $R'=\sqrt{\gamma^2R_\parallel^2+R_\perp^2}$。从静止系 $\mathbf E'=q\mathbf R'/(4\pi\varepsilon_0R'^3)$ 变换得 $E_\parallel=E_\parallel'$、$E_\perp=\gamma E_\perp'$，所以 $\mathbf E=\frac{q\gamma}{4\pi\varepsilon_0}(R_\parallel,\mathbf R_\perp)/(\gamma^2R_\parallel^2+R_\perp^2)^{3/2}$；$\mathbf B=\mathbf v\times\mathbf E/c^2$。

(2) 令 $\vartheta$ 为 $\mathbf R$ 与运动方向夹角，有 $|\mathbf E|=\frac{|q|}{4\pi\varepsilon_0R^2}\frac{1-\beta^2}{(1-\beta^2\sin^2\vartheta)^{3/2}}$。轴向强度 $E_{\rm axial}=E_C/\gamma^2$，横向 $E_{\rm transverse}=\gamma E_C$。高速时场线角向集中于横向，不能简单把全场均乘同一个 $\gamma$。

(3) 因 $\mathbf B=\mathbf v\times\mathbf E/c^2$，$\mathbf E\cdot\mathbf B=0$。又 $E^2-c^2B^2=E_\parallel^2+(1-\beta^2)E_\perp^2=E_\parallel'^2+E_\perp'^2=q^2/(16\pi^2\varepsilon_0^2R'^4)$，等于静止系电场平方，是Lorentz不变量。

(4) $\beta\to0$ 时 $\gamma\to1$，恢复Coulomb场、$\mathbf B\to0$。$\beta\to1^-$ 时在固定非横向角度下场趋零、横向增强，角分布越来越窄。此为匀速运动电荷的速度场，空间衰减均为 $R^{-2}$；加速电荷的辐射场需额外的 $R^{-1}$ 项。
